% Recovered verbatim from chapters_11_15_with_part/chapters_11_15_with_part.tex
% Canonical chapter 13
\chapter{White-State Cores in Cosmological Bounces}

\section{Bounce Mechanism}

In a contracting phase, the total energy density
\[
  \rho_{\mathrm{tot}} = \rho_\tau + \rho_\chi + \rho_{ER} + \rho_{\mathrm{m}}
\]
approaches a finite critical value $\rho_{\mathrm{crit}}$ as the white-state is approached.
Rather than diverging, geometric and chromatic terms conspire to produce a minimal scale factor $a_{\mathrm{min}}$ at which
\[
  \dot{a} = 0, \qquad \ddot{a} > 0,
\]
implementing a bounce.

\section{White-State as Nonsingular Attractor}

The white-state core $\Wcore$ acts as a nonsingular attractor configuration in the space of cosmological solutions.
Nearby trajectories in the phase--chromatic--ER configuration space are drawn into a neighbourhood of $\Wcore$ before re-expanding along new directions in field space after the bounce.

\section{Entropy and Information Across the Bounce}

Information associated with pre-bounce structures is not necessarily erased.
Instead, it can be encoded in:

\begin{itemize}
  \item the ER network pattern at the bounce,
  \item residual phase correlations,
  \item the distribution of near-white chromatic perturbations.
\end{itemize}

Thus MK--43 Ultra allows for partial memory between cycles without requiring a singularity.

